% problem-000083 7 变压吸附制氧能效
\section*{7. 变压吸附制氧能效}
⽬前家⽤⼩型制氧机⼤多以变压吸附⽅式⼯作，其原理是以空⽓为原料，在加压时吸附剂将空⽓中⼤部分氮⽓吸附，氧⽓和未被吸附的少量氮⽓被收集起来，得到富氧空⽓。被吸附的氮⽓在减压条件下通过尾⽓管排⾄⼤⽓。某⼩型家⽤常压制氧机的标称参数为额定功率100W，输出富氧空⽓中氧⽓浓度91%（体积分数），输出富氧空⽓流量为 $1 . 8 \ : \mathrm { L } \ : \mathrm { m i n ^ { - 1 } }$ ，⼯作温度为 $2 5 ~ ^ { \circ } \mathrm { C } _ { \circ }$ 以空⽓进⽓量10mol为对象，计算达到分离效果所需要的最⼩能量以及该家⽤制氧机的能量效率。空⽓由体积分数为20%的氧⽓和80%的氮⽓组成，排出的尾⽓中只有氮⽓。所有⽓体均视作理想⽓体。
